\paragraph{C.3w (Proved in the stated algebraizable scope): proper local Picard images and actual families via alterations.}
Let $k$ be algebraically closed. Let $(R,\mathfrak m,k)$
be the strict henselization at a $k$-rational point of
a normal local scheme essentially of finite type over
$k$. Assume $N=\dim R\ge3$ and that its puncture
$X=\operatorname{Spec}R\setminus\{\mathfrak m\}$ is
regular. Write
\[
                   G=\mathbf{Pic}^{\mathrm{loc}}_{R/k}
\]
for the henselized-tensor-product local Picard scheme
of C.3n. Then the following point-realization assertion holds:
\begin{equation}\label{eq:C3-proper-local-image-ordinary}
 B(k)\subset\operatorname{Pic}(X)_{\mathrm{alg}}
              \quad\text{for every closed abelian subvariety }B\subset G.
\end{equation}
This asserts ordinary algebraic equivalence of each
point, not the existence of a universal actual line
bundle on $X\times B$ representing the inclusion.

Consequently, if $\operatorname{char}k=p>0$,
C.3q closes the ordinary-surjectivity step:
for $N\ge4$,
$0\ne f\in\mathfrak m$, and regular
$Y=X\cap V(f)$, restriction
\begin{equation}\label{eq:C3-charp-ordinary-isomorphism}
 \operatorname{Pic}(X)_{\mathrm{alg}}
       \longrightarrow\operatorname{Pic}(Y)_{\mathrm{alg}}
\end{equation}
is surjective, with $p$-primary torsion kernel.
In particular $J^1(X)\to J^1(Y)$ is a rational
isomorphism. The algebraizability and specified
coefficient/residue field hypotheses in this
assertion are retained.

\begin{proof}[Proof of the point-realization assertion]
\emph{1. Fixing one regular alteration.}
Choose a normal affine algebraic model near the
specified point. The exact geometric input is
\citet[Theorem 4.1 and Remark 4.2, p.~66]{D_DeJong1996Alterations}:
it has a proper generically finite alteration
whose source is regular and quasiprojective.
Base change to $R$ preserves regularity, since
strict henselization is a localization and
ind-etale base change. Choose an integral
connected component dominating
$\operatorname{Spec}R$. We obtain
\[
             h:Z\longrightarrow\operatorname{Spec}R,
             \qquad d=[k(Z):\operatorname{Frac}R]>0,
\]
with $Z$ regular, integral, and proper over $R$.
No assertion that this alteration is finite or
is an isomorphism over the puncture is made.
Put $Z^\circ=h^{-1}(X)$ and $E=h^{-1}(\mathfrak m)$.

The Stein algebra
\[
                         R'=\Gamma(Z,\mathcal O_Z)
\]
is finite over $R$, is its normalization in
$k(Z)$, and is a normal domain. Indeed an
element integral over $R$ is regular in every
normal local ring of $Z$, and proper coherence
makes the global algebra finite. Since $R$ is
henselian, the finite domain $R'$ is local;
its residue field is $k$. Stein factorization
gives a proper birational map
$Z\to\operatorname{Spec}R'$ with connected
fibres. We use
\citet[Theorem 37.53.4]{D_StacksAlterationFamilies}
for this factorization. In particular $E$ is
connected. Choose a $k$-point $e\in E$.
Let $E_1,\ldots,E_s$ denote the prime divisors
of $Z$ contained in $E$; components of $E$
of larger codimension do not enter this list.
Each $E_i$ is Cartier because $Z$ is regular.

\emph{2. Extending a proper parameter curve as a
Picard-functor class.}
We use exactly the etale sheafification of
henselized-tensor-product Picard groups in
\citet[Introduction, Theorem 1 and preceding
definition, p.~VII]{D_Boutot1978}, as specified
in C.3n. In particular its sections have
representatives etale locally on the parameter.
Let $C$ be a smooth projective connected
$k$-curve, and let $g:C\to G$ be a morphism.
For an affine etale $U=\operatorname{Spec}B\to C$,
set
\[
 T_B=(R\otimes_kB)^h_{\mathfrak m},\quad
 Z_B^h=Z\times_R\operatorname{Spec}T_B,
 \quad (Z_B^h)^\circ=Z_B^h\setminus h^{-1}V(\mathfrak mT_B).
\]
All these rings and schemes are noetherian.
The scheme $Z_B^h$ is regular: $B/k$ is smooth,
so $R\to R\otimes_kB$ is smooth, and the
subsequent henselization is ind-etale.
Let $\mathcal F$ and $\mathcal H$ be the etale
sheaves on $C$ obtained from
\[
 U\longmapsto\operatorname{Pic}(Z_B^h),\qquad
 U\longmapsto\operatorname{Pic}((Z_B^h)^\circ).
\]
Pullback along $h$ maps the local Picard family
$g$ to a section $g'\in\mathcal H(C)$.

There is an exact sequence of sheaves on
$C_{\mathrm{et}}$
\begin{equation}\label{eq:C3-exceptional-lattice-sheaves}
 0\longrightarrow\underline{\mathbb Z}^{s}
       \longrightarrow\mathcal F\longrightarrow\mathcal H
       \longrightarrow0,
\end{equation}
where the first map sends a tuple to
$\mathcal O(\sum_i a_iE_i)$.
Surjectivity follows by closing the divisor
of a rational section on the regular scheme
$Z_B^h$, component by component.
The kernel consists of Cartier divisors
supported on the boundary. Over each connected
component of $U$, those prime boundary divisors
are exactly $E_i\times_k U$, so their coefficients
are locally constant integers.
There are no relations in this lattice.
To see this, flat base change identifies the
Stein algebra of $Z_B^h$ with
$R'\otimes_RT_B$, which is normal.
The boundary in this normal affine base has
codimension $N\ge2$. A principal divisor on
$Z_B^h$ supported on that boundary comes from
a rational function on this Stein base.
The proper birational map is an isomorphism
at its codimension-one points; thus that
function and its inverse are regular away
from a subset of codimension at least two.
Normal Hartogs extension makes it a global
unit, so its divisor is zero. This proves
injectivity in \eqref{eq:C3-exceptional-lattice-sheaves}.

One has $H^1_{\mathrm{et}}(C,\underline{\mathbb Z}^{s})=0$.
For completeness, if $j:\operatorname{Spec}k(C)\to C$
is the generic point, normality of every etale
curve over $C$ gives
$j_*\underline{\mathbb Z}^{s}=\underline{\mathbb Z}^{s}$:
its connected components and generic components
coincide. The first Leray edge map therefore
injects this $H^1$ into
\[
 H^1(k(C),\mathbb Z^s)
   =\operatorname{Hom}_{\mathrm{cont}}
              (\operatorname{Gal}(k(C)^{\mathrm{sep}}/k(C)),\mathbb Z^s)=0.
\]
A continuous homomorphism from a profinite
group to the discrete torsion-free group
$\mathbb Z^s$ has finite and hence zero image.
The exact cohomology sequence of
\eqref{eq:C3-exceptional-lattice-sheaves} now
lifts $g'$ to a section $a\in\mathcal F(C)$.
This is still a sheaf-valued Picard class;
the next step constructs its actual family.

\emph{3. Formal units on a curve have the required
cohomology after rigidification.}
Write $\widehat R'$ for the completion of $R'$;
the $\mathfrak mR'$-adic and maximal-ideal-adic
topologies agree. Let $\mathfrak n$ be its
maximal ideal. On $C_{\mathrm{et}}$ consider
the sheaves
\[
 \mathcal R=\varprojlim_r
       (\widehat R'/\mathfrak n^r)\otimes_k\mathcal O_C,
 \qquad \mathcal K=\ker(\mathcal R^\times\to\mathcal O_C^\times).
\]
Thus $\mathcal K=1+\mathfrak n\mathcal R$,
with complete separated filtration
$\mathcal K_r=1+\mathfrak n^r\mathcal R$ and
\begin{equation}\label{eq:C3-formal-unit-quotients}
 \mathcal K_r/\mathcal K_{r+1}
       = (\mathfrak n^r/\mathfrak n^{r+1})\otimes_k\mathcal O_C
\end{equation}
as additive sheaves. This description is
valid in every characteristic and uses no
logarithm.

For every affine etale $U\to C$,
\begin{equation}\label{eq:C3-formal-unit-affine-acyclicity}
                      H^i_{\mathrm{et}}(U,\mathcal K)=0
                                \quad(i>0).
\end{equation}
Here is an explicit proof avoiding a passage
from continuous cohomology to ordinary
cohomology. Take any finite affine etale
cover of $U$. Its disjoint union is affine
and faithfully flat over $U$. The augmented
Amitsur complex for each additive quotient
in \eqref{eq:C3-formal-unit-quotients} is
exact, by faithfully flat descent of modules.
Given a positive-degree Cech cocycle with
values in $\mathcal K$, solve the coboundary
equation modulo $\mathcal K_2$, lift the
cochain to $\mathcal K_1$, and divide by
its coboundary. The resulting cocycle lies
in $\mathcal K_2$. Repeat at every level.
All cochains lift on these affine intersections,
since the quotient maps of the coefficient
vector spaces split over $k$. The infinite
product of the correcting cochains converges
in the complete filtration and solves the
original coboundary equation on this same
cover. Thus every such cover has zero
positive Cech cohomology. Finite affine
etale covers form a cofinal system, stable
under fibre products. The exact criterion
\citet[Lemma 21.10.9]{D_StacksAlterationFamilies}
proves \eqref{eq:C3-formal-unit-affine-acyclicity}.

Cover $C$ by the complements of two distinct
$k$-points. Both opens and their intersection
are affine. Mayer--Vietoris and
\eqref{eq:C3-formal-unit-affine-acyclicity} give
\begin{equation}\label{eq:C3-formal-unit-curve-H2}
                           H^2_{\mathrm{et}}(C,\mathcal K)=0.
\end{equation}
In particular a gerbe banded by $\mathcal K$
on $C$ has an object: the affine etale covers
and their intersections satisfy the $H^1$
conditions, and
\citet[Lemma 21.11.1, Tag 0CK0]{D_StacksAlterationFamilies}
applies. This is the precise gerbe-neutralization
statement used below.

\emph{4. Constructing the actual formal family.}
Complete the local representatives of $a$
along $\mathfrak m$. They give a locally
represented Picard class on the formal
completion of $Z_{\widehat R}\times_k C$
along its closed fibre. Consider the stack
of its invertible-sheaf representatives,
rigidified along the closed subscheme
$\{e\}\times C$. It is locally nonempty:
after obtaining a local representative,
tensor it by the inverse of its restriction
to $e\times U$, pulled back from $U$.
That tensor factor is zero in the sheafified
Picard functor. Any two such representatives
are locally isomorphic. An isomorphism can
be made compatible with their rigidifications
by multiplying it by the inverse of its
restriction at $e$; the coefficient field
$k$ supplies the lift of this base unit.
Thus this stack is a gerbe.

Its band is exactly $\mathcal K$ of step 3.
Indeed flat base change gives
\[
 \Gamma(Z_B^h,\mathcal O)=R'\otimes_RT_B,
\]
by \citet[Lemma 30.5.2, Tag 02KH]{D_StacksAlterationFamilies}.
The theorem on formal functions then gives
the full formal global-functions algebra
\[
 \Gamma(\widehat{Z_B^h},\mathcal O)
    =\varprojlim_r (R'/\mathfrak m^rR')\otimes_kB,
\]
using \citet[Theorem 30.20.5, Tag 02OC]{D_StacksAlterationFamilies}.
This is $\mathcal R(U)$, with the equivalent
$\mathfrak n$-adic filtration. Evaluation at
$e$ is its residue map to $B$, so the
rigidification-preserving automorphisms are
$\mathcal K(U)$.
This argument identifies the inverse limit
of global functions. It does \emph{not}
claim that global functions on each individual
thickening are $R'/\mathfrak m^rR'$; such
a claim could fail through higher-cohomology
base-change terms.

Step 3 therefore neutralizes this gerbe,
producing a compatible actual formal line
bundle $\widehat{\mathcal M}$ on
$\widehat{Z_{\widehat R}\times C}$.
The scheme $Z_{\widehat R}\times C$ is
proper over the complete noetherian ring
$\widehat R$. Grothendieck existence,
\citet[Proposition 30.25.4, Tag 088C]{D_StacksAlterationFamilies},
algebraizes $\widehat{\mathcal M}$ to an
invertible sheaf $\mathcal M$ on
$Z_{\widehat R}\times C$.
Invertibility follows either by algebraizing
its formal inverse and the two evaluation
isomorphisms, or by properness and the
closed non-locally-free locus.
For every $t\in C(k)$, its restriction to
$Z^\circ_{\widehat R}\times\{t\}$ represents
the pullback of $g(t)\in G(k)=\operatorname{Pic}(X)$
after completion. This follows from the
definition of the lifted class $a$; at a
$k$-point etale sheafification adds no new
isomorphism classes.

\emph{5. Rigidity of proper-curve differences
under approximation.}
We need a precise fact about a proper scheme
$W$ over a complete noetherian local
$k$-algebra with residue $k$ and connected
closed fibre. Suppose two line bundles
$\mathcal L_1,\mathcal L_2$ on $W\times C$
agree on the closed fibre times $C$.
Then for every $t_0,t_1\in C(k)$ their
fibre differences agree in $\operatorname{Pic}(W)$:
\begin{equation}\label{eq:C3-proper-curve-rigidity}
 [\mathcal L_{1,t_1}\otimes\mathcal L_{1,t_0}^{-1}]
   =[\mathcal L_{2,t_1}\otimes\mathcal L_{2,t_0}^{-1}].
\end{equation}

To prove this, let $W_r$ be the closed
infinitesimal thickenings and put
$\mathcal D=\mathcal L_1\otimes\mathcal L_2^{-1}$.
Each $W_r$ is a proper finite-type $k$-scheme,
so its Picard functor is a separated group
scheme locally of finite type
\citep[Proposition 1.3, p.~4]{D_ChengJiLarsonOlander2021}.
The kernel of
$\operatorname{Pic}_{W_r/k}\to\operatorname{Pic}_{W_1/k}$
is quasi-affine of finite presentation.
For this precise property apply
\citet[Proposition 1.4, p.~4, from Theorem 1.2,
p.~3]{D_ChengJiLarsonOlander2021}
to the proper surjective closed immersion
$W_1\hookrightarrow W_r$: its base is
$\operatorname{Spec}k$, so the theorem's
nonempty base open is the whole base.
No smoothness or reducedness of the
thickenings is required.

The family $\mathcal D$ defines a morphism
$C\to\operatorname{Pic}_{W_r/k}$ whose
image is in this quasi-affine kernel.
Every morphism from the proper connected
curve $C$ to a quasi-affine $k$-scheme is
constant: compose with an open immersion
into an affine scheme and use
$\Gamma(C,\mathcal O_C)=k$.
Thus $\mathcal D_{t_1}\otimes\mathcal D_{t_0}^{-1}$
is trivial on every $W_r$.
These trivializations can be chosen compatibly.
For fixed $r$, the images of
$\Gamma(W_m,\mathcal O)\to\Gamma(W_r,\mathcal O)$
stabilize, since they are descending
subspaces of a finite-dimensional $k$-vector
space. These global-function algebras are
local Artinian with residue $k$, because
the common underlying closed fibre is
connected. Hence images of their unit
groups stabilize as well: an element of
an image algebra is a unit precisely when
its residue is nonzero, and then any
preimage is a unit. The nonempty torsors
of trivializations therefore satisfy the
Mittag--Leffler condition and admit a
compatible inverse-limit point.
Formal full faithfulness
\citep[Lemma 30.24.1, Tag 0883]{D_StacksFormalPicard}
now proves \eqref{eq:C3-proper-curve-rigidity}.

The line bundle $\mathcal M$ of step 4
descends, with its invertibility, to
$Z_B\times C$ for a finitely generated
$R$-subalgebra $B\subset\widehat R$:
this is finite-presentation descent for
line bundles on the fixed quasi-compact
quasi-separated scheme $Z\times C$.
Since $R$ is an excellent henselian local
ring, Artin approximation in the precise
form
\citet[Theorem 16.13.1, Tag 07QY]{D_StacksAlterationFamilies}
supplies an $R$-map $B\to R$ agreeing
modulo $\mathfrak m$ with $B\to\widehat R$.
Pulling back gives an actual invertible
sheaf $\mathcal M_R$ on $Z\times C$
whose completion agrees with $\mathcal M$
on the closed fibre times $C$.
Apply \eqref{eq:C3-proper-curve-rigidity}
to $W=Z_{\widehat R}$. The completed fibre
differences of $\mathcal M_R$ are exactly
those of $\mathcal M$, hence on
$Z^\circ_{\widehat R}$ are the pullbacks
of $g(t_1)-g(t_0)$ after completion.

\emph{6. Transfer of the actual family.}
There is an open $U\subset X$ containing
all points of codimension at most one
such that $h^{-1}(U)\to U$ is finite flat
of degree $d$. Indeed, under the finite
normalization over the normal domain $R$,
primes above a codimension-one prime have
codimension one. Their local rings in
$R'$ are discrete valuation rings, and
the proper birational map
$Z\to\operatorname{Spec}R'$ is an
isomorphism there. Over that codimension-one
point of the regular $X$, the finite
normalization algebra is torsion-free
over a discrete valuation ring and is
therefore flat. Properness makes
the finite locus open, and flatness then
gives the required $U$.

On the actual product $U\times C$, form
the finite-flat norm of
$\mathcal M_R|_{h^{-1}(U)\times C}$.
Explicitly, for a finite flat map $v$,
\[
                 \operatorname{Nm}_v(\mathcal L)
                  =\det(v_*\mathcal L)\otimes\det(v_*\mathcal O)^{-1}.
\]
This norm commutes with arbitrary base
change and satisfies
$\operatorname{Nm}_v(v^*M)=M^d$.
The complement of $U\times C$ has
codimension at least two in the regular
scheme $X\times C$. Closure of divisors
therefore extends the norm uniquely to
an invertible sheaf $\mathcal N$ on
$X\times C$. The same uniqueness shows
that its fibres are the corresponding
extensions of the fibre norms.

Completion is flat and regular here.
The complement of $U_{\widehat R}$ still
has codimension at least two, by going
down, and $X_{\widehat R}$ is regular.
Thus the extended norm also commutes
with completion: this can be checked on
the finite-flat open and extended
uniquely. Step 5 and the norm identity
therefore give
\[
 [\mathcal N_{t_1}\otimes\mathcal N_{t_0}^{-1}]_{\widehat R}
                  =d\bigl(g(t_1)-g(t_0)\bigr)_{\widehat R}.
\]
Elkik's Picard-group comparison for a
noetherian henselian pair,
\citet[Section III, unnumbered corollary,
p.~573]{D_Elkik1973}, makes
$\operatorname{Pic}(X)\to\operatorname{Pic}(X_{\widehat R})$
injective. Consequently
\begin{equation}\label{eq:C3-alteration-curve-multiple}
 d\bigl(g(t_1)-g(t_0)\bigr)
      =[\mathcal N_{t_1}\otimes\mathcal N_{t_0}^{-1}]
             \in\operatorname{Pic}(X)_{\mathrm{alg}}.
\end{equation}
Every line bundle used in this last
ordinary relation is on the fixed actual
product $X\times C$.

\emph{7. Removing the degree and applying C.3q.}
Let $b\in B(k)$ for a closed abelian
subvariety $B\subset G$. Multiplication
by $d$ on $B(k)$ is surjective, including
when $d$ is divisible by the characteristic.
Choose $b'\in B(k)$ with $db'=b$.
Join $0$ and $b'$ by a chain of integral
curves in the projective variety $B$,
and normalize them. These are smooth
projective curves over the perfect field
$k$. Apply \eqref{eq:C3-alteration-curve-multiple}
to their maps to $G$ and add the
resulting actual ordinary relations.
Their sum is $db'=b$. This proves
\eqref{eq:C3-proper-local-image-ordinary}.

Finally suppose $\operatorname{char}k=p>0$,
$N\ge4$, and $Y$ is regular. C.3n places
every ordinary target class in a closed
abelian subvariety of the target local
Picard scheme. C.3q gives a surjective
isogeny from a closed source abelian
subvariety onto that target subvariety.
Every $k$-point lifts under this isogeny,
and \eqref{eq:C3-proper-local-image-ordinary}
makes the lifted point ordinary on $X$.
Thus \eqref{eq:C3-charp-ordinary-isomorphism}
is surjective integrally. C.3n identifies
its kernel as $p$-primary torsion, giving
the rational isomorphism asserted above.
\end{proof}

\emph{Scope of the construction.}
The proof uses an actual regular
alteration of an algebraic model, a
constant exceptional-divisor lattice,
neutralization of a specifically
rigidified formal Picard gerbe, proper
formal existence, and Artin approximation
of a whole proper-curve family. None of
these steps is a blanket essential
surjectivity assertion for henselian
GAGA. The transfer is taken only after
the actual family on $Z\times C$ has
been constructed. No all-excellent-ring
extension beyond the stated model is
claimed here.
